\resizebox{0.96\textwidth}{!}{%
\begin{tikzpicture}[
    scale=0.94,
    sectiontitle/.style={anchor=west, font=\bfseries\normalsize, text=black},
    actor/.style={draw=black, fill=gray!8, thick, rounded corners=5pt, minimum width=3.2cm, minimum height=1.1cm, align=center, font=\small, text=black},
    eve/.style={draw=black, fill=gray!22, very thick, double, rounded corners=5pt, minimum width=4.0cm, minimum height=1.1cm, align=center, font=\small, text=black},
    note/.style={draw=black!75, fill=white, font=\scriptsize, align=center, inner sep=5pt, rounded corners=3pt, text=black},
    evenote/.style={draw=black, fill=gray!12, dashed, font=\scriptsize, align=center, inner sep=5pt, rounded corners=3pt, text=black},
    warnnote/.style={draw=black, fill=gray!18, thick, font=\scriptsize, align=center, inner sep=5pt, rounded corners=4pt, text=black},
    msg/.style={-{Stealth[scale=1.15]}, thick, black},
    evemsg/.style={-{Stealth[scale=1.15]}, very thick, dashed, black},
    lifeline/.style={thick, dotted, black!50},
    evelifeline/.style={thick, dashed, black!65}
]

    % ==============================================================
    % (a) KỊCH BẢN BÌNH THƯỜNG (LÝ TƯỞNG)
    % ==============================================================
    \node[sectiontitle] at (0, 13.0) {(a) Trao đổi bình thường qua kênh truyền không bảo vệ (Kịch bản lý tưởng)};

    % Diễn viên
    \node[actor] (bobA) at (2.2, 11.6) {\textbf{Bob}\\[2pt]{\scriptsize Cặp khóa: $(PU_B, PR_B)$}};
    \node[actor] (aliceA) at (12.3, 11.6) {\textbf{Alice}\\[2pt]{\scriptsize Nhận đúng khóa $PU_B$}};

    % Đường thời gian (Lifelines)
    \draw[lifeline] (2.2, 11.0) -- (2.2, 7.9);
    \draw[lifeline] (12.3, 11.0) -- (12.3, 7.9);

    % Bước 1: Bob gửi PU_B
    \draw[msg] (2.2, 9.9) -- (12.3, 9.9)
        node[midway, above=3pt, font=\footnotesize] {1. Bob gửi khóa công khai $PU_B$};

    % Bước 2: Alice gửi bản mã C
    \draw[msg] (12.3, 8.6) -- (2.2, 8.6)
        node[midway, above=3pt, font=\footnotesize] {2. Alice gửi bản mã $C = E(PU_B, M)$};

    % Ghi chú giải mã & xác thực
    \node[note, text width=4.0cm] at (2.2, 7.3) {
        \textbf{Bob giải mã bằng $PR_B$:}\\
        $M = D(PR_B, C)$
    };
    \node[note, text width=4.0cm] at (12.3, 7.3) {
        \textbf{Alice an tâm:}\\
        Chỉ người có $PR_B$ mới giải mã được
    };

    % Đường phân cách hai kịch bản
    \draw[dash pattern=on 5pt off 4pt, black!50, thick] (0, 6.1) -- (14.5, 6.1);

    % ==============================================================
    % (b) KỊCH BẢN TẤN CÔNG MAN-IN-THE-MIDDLE (EVE TRÁO ĐỔI KHÓA)
    % ==============================================================
    \node[sectiontitle] at (0, 5.3) {(b) Kịch bản tấn công Man-In-The-Middle (Eve can thiệp và tráo đổi khóa)};

    % Diễn viên
    \node[actor] (bobB) at (1.8, 3.9) {\textbf{Bob}\\[2pt]{\scriptsize Cặp khóa: $(PU_B, PR_B)$}};
    \node[eve]   (eveB) at (7.25, 3.9) {\textbf{Eve (Kẻ tấn công MITM)}\\[2pt]{\scriptsize Cặp khóa giả mạo: $(PU_E, PR_E)$}};
    \node[actor] (aliceB) at (12.7, 3.9) {\textbf{Alice}\\[2pt]{\scriptsize Nhầm $PU_E$ là khóa của Bob}};

    % Đường thời gian (Lifelines)
    \draw[lifeline] (1.8, 3.3) -- (1.8, -4.5);
    \draw[evelifeline] (7.25, 3.3) -- (7.25, 2.85);
    \draw[evelifeline] (7.25, 2.15) -- (7.25, -0.95);
    \draw[lifeline] (12.7, 3.3) -- (12.7, -4.5);

    % Bước 1: Bob gửi PU_B -> Eve chặn
    \node[evenote, text width=2.6cm] (chaneve) at (7.25, 2.5) {\textbf{Eve chặn bắt $PU_B$}};
    \draw[msg] (1.8, 2.5) -- (chaneve.west)
        node[midway, above=3pt, font=\footnotesize] {1. Bob gửi khóa $PU_B$};

    % Bước 2: Eve thay bằng PU_E gửi Alice
    \draw[evemsg] (7.25, 1.1) -- (12.7, 1.1)
        node[midway, above=3pt, font=\footnotesize] {2. Gửi $PU_E$ (mạo danh Bob)};

    % Bước 3: Alice gửi C1 về cho Eve (nhầm là Bob)
    \draw[msg] (12.7, -0.3) -- (7.25, -0.3)
        node[midway, above=3pt, font=\footnotesize, fill=white, inner sep=2pt] {3. Gửi bản mã $C_1 = E(PU_E, M)$};

    % Hộp Eve giải mã, đọc trộm và mã hóa lại
    \node[evenote, text width=6.6cm] (evedecrypt) at (7.25, -1.8) {
        \textbf{Eve giải mã bằng $PR_E$ \& đọc trộm:}\\
        $M = D(PR_E, C_1)$ \textit{(Thu được thông điệp gốc $M$)}\\[3pt]
        \textbf{Eve mã hóa lại bằng $PU_B$ thật của Bob:}\\
        $C_2 = E(PU_B, M)$
    };
    \draw[evelifeline] (evedecrypt.south) -- (7.25, -4.5);

    % Bước 4: Eve gửi C2 cho Bob
    \draw[evemsg] (7.25, -3.6) -- (1.8, -3.6)
        node[midway, above=3pt, font=\footnotesize, fill=white, inner sep=2pt] {4. Chuyển tiếp $C_2 = E(PU_B, M)$};

    % Hàng kết quả cuối cùng: 3 khối cân xứng
    \node[note, text width=3.2cm] at (1.8, -5.3) {
        \textbf{Bob giải mã ($PR_B$):}\\
        $M = D(PR_B, C_2)$\\
        {\scriptsize (Nhận đúng $M$, không nghi ngờ)}
    };

    \node[warnnote, text width=6.2cm] at (7.25, -5.3) {
        \textbf{Hậu quả an ninh nghiêm trọng:}\\
        Eve đọc/sửa trộm $M$ mà hai bên hoàn toàn không biết\\
        do thiếu cơ chế xác thực danh tính khóa công khai!
    };

    \node[note, text width=3.2cm] at (12.7, -5.3) {
        \textbf{Alice tin tưởng:}\\
        Tin Bob nhận được $C_1$\\
        {\scriptsize (Không biết bị tráo khóa $PU_E$)}
    };

\end{tikzpicture}%
}
